% C6_brief.tex -- ONE source for the English and Romanian case brief of Case 6 (Lecture 6).
\input{numbers_\lang}
\renewcommand{\emqacase}{\T{Research case 6 (Lecture 6)}{Studiul de caz 6 (cursul 6)}}
\begin{document}
\emqatitle{\T{Case 6: Should we buy a battery? Probabilistic forecasts and day-ahead arbitrage}%
                {Studiul de caz 6: Merită investiția într-o baterie? Prognoze probabilistice și arbitraj pe piața pentru ziua următoare}}%
          {\T{Based on}{Studiu bazat pe lucrarea} Andrei, Cepoi, Dumitrescu, Lessmann, Chiriac (Luchian), Pele \& Petukhina (2026), \emph{Joint probabilistic forecasting for battery arbitrage in day-ahead power markets},
           SSRN 6872381, \href{https://doi.org/10.2139/ssrn.6872381}{doi:10.2139/ssrn.6872381}}

\begin{scenario}
\T{A Romanian producer with a solar park is offered a battery energy storage system (BESS). The vendor's slide shows the profit a battery \emph{could} have made in 2022.
The investment committee must decide: \textbf{how much does a battery really earn from buying cheap and selling dear on the day-ahead market (PZU), when it must bid a day ahead without knowing tomorrow's prices --- and does that repay the investment?}
The trading desk asks a second question: \textbf{which hours should the battery bid, and at what prices?}}%
{Unui producător din România care deține un parc fotovoltaic i se propune un sistem de stocare a energiei în baterii (BESS, \emph{battery energy storage system}). Prezentarea ofertantului arată profitul pe care bateria l-\emph{ar fi putut} obține în 2022.
Comitetul de investiții trebuie să decidă: \textbf{cît cîștigă efectiv o baterie care cumpără ieftin și vinde scump pe piața pentru ziua următoare (PZU), dacă ofertele se depun cu o zi înainte, fără a cunoaște prețurile de a doua zi --- și se recuperează astfel investiția?}
Departamentul de tranzacționare are o a doua întrebare: \textbf{în ce ore și la ce prețuri ar trebui să liciteze bateria?}}
\end{scenario}

\section{\T{What the paper did}{Lucrarea de referință}}
\textbf{\T{Question.}{Întrebarea de cercetare.}} \T{How much arbitrage value can be captured \emph{ex ante} when storage is operated on probabilistic price forecasts, and does it justify the investment, in four interconnected Central and Eastern European (CEE) day-ahead markets (Sections~1 and 3)?}%
{Ce parte din valoarea de arbitraj poate fi captată \emph{ex ante} atunci cînd sistemul de stocare este exploatat pe baza prognozelor probabilistice ale prețului și justifică aceasta investiția, pe patru piețe interconectate pentru ziua următoare din Europa Centrală și de Est (secțiunile~1 și 3)?}

\textbf{\T{Data.}{Datele.}} \T{Hourly day-ahead prices of Romania, Bulgaria, Hungary and the Czech Republic, January 2018 -- September 2024 (59,160 hours per country), with 37 predictors: load, generation and their forecasts, lagged prices, Brent, TTF and carbon prices, weather and calendar (Section~5.1).
Battery: 1 MW / 1 MWh, efficiency 90\,\%, CAPEX (capital expenditure) 500,000 EUR, 15 years, 7\,\% discount rate, at most one cycle per day (Section~4).}%
{Prețurile orare pentru ziua următoare din România, Bulgaria, Ungaria și Cehia, ianuarie 2018 -- septembrie 2024 (59.160 de ore pentru fiecare țară), cu 37 de variabile explicative: consumul, producția și prognozele acestora, prețurile anterioare, prețurile Brent, TTF și ale certificatelor de emisii de CO$_2$, variabile meteorologice și calendaristice (secțiunea~5.1).
Bateria: 1 MW / 1 MWh, randament de 90\,\%, CAPEX (cheltuieli de capital) de 500.000 EUR, durată de viață de 15 ani, rată de actualizare de 7\,\%, cel mult un ciclu pe zi (secțiunea~4).}

\textbf{\T{Method.}{Metoda.}} \T{Distributional Random Forests (DRF) forecast the \emph{joint} distribution of the 96 next-day prices (24 hours $\times$ 4 countries), refitted on the last 365 days every 30 days.
A quantile-based (QB) rule chooses the buy hour and the sell hour from the forecast medians and places limit bids at forecast quantiles; benchmarks are a fixed pair of hours (NH) and daily perfect foresight (PF-Daily).}%
{Modelele DRF (\emph{Distributional Random Forests}, păduri aleatoare care estimează distribuții întregi) prognozează distribuția \emph{comună} a celor 96 de prețuri din ziua următoare (24 de ore $\times$ 4 țări) și sînt reestimate la fiecare 30 de zile pe ultimele 365 de zile.
O regulă bazată pe cuantile (QB) alege ora de cumpărare și ora de vînzare pe baza medianelor prognozate și plasează oferte cu preț limită la cuantilele prognozei; reperele de comparație sînt o pereche fixă de ore (NH) și previziunea perfectă zilnică (PF-Daily).}

\textbf{\T{Main findings of the paper.}{Principalele rezultate ale lucrării.}}
\begin{itemize}
\item \T{``At $\alpha=0.8$, QB beats NH in cumulative profit in all four markets, with relative gains of approximately 13\,\% in Czechia \dots{} and 7\,\% in Romania'' (Section~6.2).
Romania, January 2019 -- September 2024: QB $\alpha=0.8$ 142,671 EUR, NH 133,230 EUR, PF-Daily 188,249 EUR; payback 20.1 years for QB (Table ``single country strategy, RO'').}%
{Pentru $\alpha=0{,}8$, strategia QB obține un profit cumulat mai mare decît NH pe toate cele patru piețe, cu cîștiguri relative de aproximativ 13\,\% în Cehia și 7\,\% în România (secțiunea~6.2).
România, ianuarie 2019 -- septembrie 2024: QB $\alpha=0{,}8$ 142.671 EUR, NH 133.230 EUR, PF-Daily 188.249 EUR; perioada de recuperare a investiției este de 20,1 ani pentru QB (tabelul pentru România).}
\item \T{``Under the baseline CAPEX of 500 thousand EUR, none of the strategies is able to cover investment costs'' (Section~6.2); break-even CAPEX lies between 48\,\% and 75\,\% of the baseline for the QB rule (Section~6.2).}%
{La CAPEX-ul de referință de 500.000 EUR, nicio strategie nu acoperă costul investiției (secțiunea~6.2); pentru regula QB, CAPEX-ul de echilibru se situează între 48\,\% și 75\,\% din valoarea de referință (secțiunea~6.2).}
\item \T{Joint cross-border dispatch adds about 5\,\% to QB profit (149.9k vs 142.7k EUR for Romania) and 18\,\% under perfect foresight, ``driven by persistent regional price differences'' (Abstract; Table ``Multi-zone joint-dispatch'').}%
{Operarea comună, transfrontalieră, crește profitul QB cu aproximativ 5\,\% (149,9 mii EUR față de 142,7 mii EUR pentru România) și cu 18\,\% în ipoteza previziunii perfecte, datorită diferențelor regionale persistente de preț (rezumat; tabelul pentru operarea în mai multe zone).}
\item \T{DRF accuracy for Romania: MAE 39.63, RMSE 49.60, CRPS 28.12 EUR/MWh over 70 rolling windows (Table ``DRF forecasting performance'').}%
{Acuratețea DRF pentru România: MAE 39,63, RMSE 49,60, CRPS 28,12 EUR/MWh, pe 70 de ferestre mobile (tabelul privind performanța prognozelor DRF).}
\end{itemize}
{\footnotesize\color{FooterGray}\T{Related work of our group for Lecture 6: Basangova, Bolovăneanu, Conda, Pele et al.\ (2026), \emph{Day-ahead forecasting for redispatch measures using machine learning models}, SSRN 7511080, \href{https://doi.org/10.2139/ssrn.7511080}{doi:10.2139/ssrn.7511080} --- neural networks (NBEATSx, NHiTS, TFT) forecast German redispatch volumes one day ahead better than earlier-generation models (Abstract; Table~4).}%
{O altă lucrare a grupului nostru relevantă pentru cursul 6: Basangova, Bolovăneanu, Conda, Pele et al.\ (2026), \emph{Day-ahead forecasting for redispatch measures using machine learning models}, SSRN 7511080, \href{https://doi.org/10.2139/ssrn.7511080}{doi:10.2139/ssrn.7511080} --- rețelele neuronale (NBEATSx, NHiTS, TFT) prognozează volumele de redispecerizare din Germania pentru ziua următoare mai bine decît modelele din generațiile anterioare (rezumat; Tabelul~4).}}

\section{\T{The method, step by step}{Metoda, pas cu pas}}
\begin{intuition}
\T{A battery earns the spread between a cheap hour and an expensive hour, minus the 10\,\% lost in charging and discharging. Looking back, the best two hours are obvious; looking forward, they are a guess.
A probabilistic forecast says not only ``the price at 19:00 will be about 150'' but ``with 80\,\% probability between 110 and 210'' --- which tells the trader how aggressively to bid.}%
{O baterie cîștigă diferența dintre prețul unei ore ieftine și cel al unei ore scumpe, minus cei 10\,\% pierduți la încărcare și descărcare. Retrospectiv, cele mai bune două ore sînt evidente; în momentul deciziei, ele sînt doar o estimare.
O prognoză probabilistică nu spune doar „prețul la ora 19:00 va fi de aproximativ 150”, ci „cu o probabilitate de 80\,\%, prețul va fi între 110 și 210” --- informație care îi arată responsabilului cu tranzacționarea cît de agresiv poate licita.}
\end{intuition}
\T{Let $p_{d,h}$ be the day-ahead price of hour $h$ on day $d$ and $\eta=0.9$ the efficiency. Buying 1 MWh in hour $h_1$ and selling in hour $h_2>h_1$ earns $\pi_d=\eta\,p_{d,h_2}-p_{d,h_1}/\eta$.
\emph{Perfect foresight} takes the best pair each day (an upper bound); \emph{fixed hours} use one pair, chosen \emph{ex ante} as the best pair on average over the previous year (the 365 days before 8 January 2019 for the 2019--2024 window, the year 2024 for 2025--2026), and trade every day.
The random forest gives each past day $i$ a weight $w_i(\mathbf x_d)$ --- the share of trees in which day $i$ falls in the same leaf as tomorrow's conditions $\mathbf x_d$ --- so the forecast is a whole distribution of 24-hour price vectors, and its quantiles $\hat Q_{\tau,d,h}$ are weighted quantiles.
The forecast-based rule picks $h_1<h_2$ maximising $\eta\hat Q_{0.5,d,h_2}-\hat Q_{0.5,d,h_1}/\eta$, trades only if this is positive, and bids to buy at $\hat Q_{1-q,d,h_1}$ and to sell at $\hat Q_{q,d,h_2}$, with $q=(1-\alpha)/2$: a larger $\alpha$ means safer bids that clear more often.
Investment: annual profit $\bar\pi_{ann}=365\,\bar\pi_d$, $\mathrm{NPV}=-\mathrm{CAPEX}+\sum_{t=1}^{15}\bar\pi_{ann}/1.07^t$, payback $=\mathrm{CAPEX}/\bar\pi_{ann}$. Forecast quality: CRPS (continuous ranked probability score) $=\sum_iw_i|y_i-y|-\tfrac12\sum_{i,k}w_iw_k|y_i-y_k|$, in EUR/MWh; lower is better.}%
{Fie $p_{d,h}$ prețul pentru ziua următoare în ora $h$ a zilei $d$ și $\eta=0{,}9$ randamentul. Cumpărarea a 1 MWh în ora $h_1$ și vînzarea lui în ora $h_2>h_1$ aduc $\pi_d=\eta\,p_{d,h_2}-p_{d,h_1}/\eta$.
Strategia de \emph{previziune perfectă} alege în fiecare zi cea mai bună pereche de ore (o limită superioară); strategia cu \emph{ore fixe} folosește o singură pereche, aleasă \emph{ex ante} ca fiind, în medie, cea mai bună în anul anterior (cele 365 de zile dinaintea datei de 8 ianuarie 2019 pentru perioada 2019--2024, anul 2024 pentru 2025--2026), și tranzacționează în fiecare zi.
Modelul Random Forest atribuie fiecărei zile trecute $i$ o pondere $w_i(\mathbf x_d)$ --- proporția arborilor în care ziua $i$ ajunge în aceeași frunză ca și condițiile din ziua următoare, $\mathbf x_d$ ---, astfel că prognoza este o distribuție completă a vectorilor de prețuri pe 24 de ore, iar cuantilele sale $\hat Q_{\tau,d,h}$ sînt cuantile ponderate.
Regula bazată pe prognoză alege $h_1<h_2$ care maximizează $\eta\hat Q_{0.5,d,h_2}-\hat Q_{0.5,d,h_1}/\eta$, tranzacționează numai dacă această valoare este pozitivă și licitează cumpărarea la $\hat Q_{1-q,d,h_1}$ și vînzarea la $\hat Q_{q,d,h_2}$, cu $q=(1-\alpha)/2$: un $\alpha$ mai mare înseamnă oferte mai prudente, acceptate mai des.
Investiția: profitul anual este $\bar\pi_{ann}=365\,\bar\pi_d$, $\mathrm{VAN}=-\mathrm{CAPEX}+\sum_{t=1}^{15}\bar\pi_{ann}/1{,}07^t$, iar perioada de recuperare $=\mathrm{CAPEX}/\bar\pi_{ann}$. Calitatea prognozei se măsoară prin CRPS (\emph{continuous ranked probability score}, scorul probabilistic ordonat continuu) $=\sum_iw_i|y_i-y|-\tfrac12\sum_{i,k}w_iw_k|y_i-y_k|$, în EUR/MWh; valorile mici sînt preferabile.}
\begin{decision}
\T{Value the battery on what a forecast-driven strategy earns, never on perfect foresight. Then compare the annual profit with the \NneedYear\ EUR per year needed to repay 500,000 EUR in 15 years at 7\,\%.}%
{Evaluați bateria după profitul unei strategii bazate pe prognoză, niciodată după cel obținut în ipoteza previziunii perfecte. Comparați apoi profitul anual cu cei \NneedYear\ EUR pe an necesari pentru a recupera 500.000 EUR în 15 ani, la o rată de 7\,\%.}
\end{decision}

\section{\T{Replicate it yourself}{Replicare pe datele cursului}}
\T{Notebook \texttt{C6\_notebook\_EN.ipynb}: hourly day-ahead prices of RO, BG and HU (energy-charts.info), CET hours; the paper's test window 8 January 2019 -- 30 September 2024 (\NnDays\ days) and an extension to October 2026.
Classroom implementation: a standard random forest with leaf weights instead of DRF; only lagged prices and calendar as predictors; no Czech data; a simplified battery state (empty/full) instead of the paper's three-level rule.}%
{Notebook-ul \texttt{C6\_notebook\_RO.ipynb} folosește prețurile orare pentru ziua următoare din România, Bulgaria și Ungaria (energy-charts.info), cu orele exprimate în CET; perioada de test din lucrare, 8 ianuarie 2019 -- 30 septembrie 2024 (\NnDays\ de zile), și o extensie pînă în octombrie 2026.
Față de lucrare, implementarea de la curs folosește un model Random Forest obișnuit, cu ponderi derivate din frunze, în locul DRF; ca variabile explicative, numai prețurile anterioare și variabilele calendaristice; fără Cehia; o stare simplificată a bateriei (goală sau plină) în locul regulii cu trei niveluri din lucrare.}

\begin{figure}[H]
\begin{minipage}[t]{0.49\linewidth}\centering
\includegraphics[width=\linewidth]{f1_daily_shape}
\caption{\T{Mean hourly price shape in Romania, three periods.}{Profilul orar mediu al prețului în România, în trei perioade.}}
\src{energy-charts.info; 2019--Oct 2026.}
\end{minipage}\hfill
\begin{minipage}[t]{0.49\linewidth}\centering
\includegraphics[width=\linewidth]{f2_forecast_fan}
\caption{\T{A probabilistic forecast of tomorrow's 24 prices.}{O prognoză probabilistică a celor 24 de prețuri din ziua următoare.}}
\src{\T{our random forest; energy-charts.info.}{modelul nostru Random Forest; energy-charts.info.}}
\end{minipage}
\end{figure}

\T{\textbf{Forecast quality.} Over the paper's window our forest has MAE \NmaeRF, RMSE \NrmseRF\ and CRPS \NcrpsRF\ EUR/MWh (paper's DRF: 39.63, 49.60, 28.12 --- different predictors and averaging), and its 10\,\%--90\,\% band covers \NcoverRF\,\% of hours (target 80\,\%).
Note: the naive ``same hour yesterday'' has MAE \NmaeNaive, so our point forecast adds little; its value is in the \emph{distribution}, which sets the bids.}%
{\textbf{Calitatea prognozei.} În perioada de test din lucrare, modelul nostru are MAE \NmaeRF, RMSE \NrmseRF\ și CRPS \NcrpsRF\ EUR/MWh (DRF în lucrare: 39,63, 49,60, 28,12 --- cu alte variabile explicative și altă metodă de mediere), iar intervalul dintre cuantilele sale de 10\,\% și 90\,\% cuprinde prețul efectiv în \NcoverRF\,\% din ore (nivelul-țintă este 80\,\%).
De reținut: prognoza naivă „aceeași oră din ziua precedentă” are MAE \NmaeNaive, deci prognoza noastră punctuală aduce un cîștig mic; valoarea ei constă în \emph{distribuție}, pe baza căreia se stabilesc ofertele.}

\begin{table}[H]\centering\small
\caption{\T{Battery strategies in Romania, 1 MW / 1 MWh, 8 Jan 2019 -- 30 Sep 2024, and annual profit in 2025--Oct 2026. Best realistic or overall value in bold.}%
{Strategiile de operare a bateriei în România, 1 MW / 1 MWh, 8 ian.\ 2019 -- 30 sept.\ 2024, și profitul anual în 2025--oct.\ 2026. Cea mai bună valoare este marcată cu aldine.}}
\tabin{strategies}
\src{\T{our computation; CAPEX 500,000 EUR, 15 years, 7\,\%. Paper (RO, same window): PF-Daily 188,249; NH 133,230; QB $\alpha$=0.8 142,671 EUR.}{calcule proprii; CAPEX 500.000 EUR, 15 ani, 7\,\%. Lucrarea (RO, aceeași perioadă): PF-Daily 188.249; NH 133.230; QB $\alpha=0{,}8$ 142.671 EUR.}}
\end{table}

\begin{figure}[H]
\begin{minipage}[t]{0.49\linewidth}\centering
\includegraphics[width=\linewidth]{f3_cumulative_profit}
\caption{\T{Cumulative profit by strategy; dotted: end of the paper's window.}{Profitul cumulat, pe strategii; linia punctată marchează sfîrșitul perioadei analizate în lucrare.}}
\src{\T{our computation; energy-charts.info, 2019--Oct 2026.}{calcule proprii; energy-charts.info, 2019--oct.\ 2026.}}
\end{minipage}\hfill
\begin{minipage}[t]{0.49\linewidth}\centering
\includegraphics[width=\linewidth]{f4_profit_by_year}
\caption{\T{Profit per year vs the annual profit needed to repay the battery.}{Profitul anual comparativ cu profitul anual necesar pentru recuperarea investiției în baterie.}}
\src{\T{our computation; 2026 = January--October.}{calcule proprii; 2026 = ianuarie--octombrie.}}
\end{minipage}
\end{figure}

\T{\textbf{What our data show.} Our perfect-foresight profit, \NpfCum\ EUR, and our fixed-hours profit, \NnhCum\ EUR (hours \NnhPair\ CET), are close to the paper's 188,249 and 133,230 EUR.
Our forecast-based rule with $\alpha=0.8$ earns \NqbCum\ EUR (\NqbVsNh\,\% vs fixed hours), below the paper's 142,671: with only past prices as predictors, our forecasts cannot beat the fixed schedule, whereas the paper's DRF with fundamentals beats it by 7\,\%.
Safe bids matter: $\alpha=0$ earns only \NqbZero\ EUR, because aggressive bids rarely clear (both bids accepted on \NqbBoth\,\% of days at $\alpha=0.8$).
Cross-border, the best of the RO/BG/HU routes adds \NcrossGain\,\% to the perfect-foresight profit; the most frequent cross-border route is \NtopRoute.
\textbf{Decision number:} 2019--2024, the realistic strategy earned \NqbYear\ EUR/year against \NneedYear\ needed (NPV \NqbNPV\ EUR, payback \NqbPay\ years) --- not viable, as in the paper.
In 2025--2026, with deeper solar middays (the fixed pair chosen on 2024 is \NnhPairNew), it earns \NqbYearNew\ EUR/year (payback \NqbPayNew\ years): closer, but still below the threshold without balancing or capacity revenues.}%
{\textbf{Ce arată datele.} Profitul nostru în ipoteza previziunii perfecte, de \NpfCum\ EUR, și profitul cu ore fixe, de \NnhCum\ EUR (orele \NnhPair\ CET), sînt apropiate de cele din lucrare: 188.249 și 133.230 EUR.
Regula noastră bazată pe prognoză, cu $\alpha=0{,}8$, cîștigă \NqbCum\ EUR (\NqbVsNh\,\% față de orele fixe), mai puțin decît cei 142.671 EUR din lucrare: avînd ca variabile explicative numai prețurile anterioare, prognozele noastre nu pot depăși programul fix, în timp ce modelul DRF din lucrare, care include variabile fundamentale, îl depășește cu 7\,\%.
Prudența ofertelor contează: pentru $\alpha=0$, strategia cîștigă numai \NqbZero\ EUR, deoarece ofertele agresive sînt rareori acceptate (pentru $\alpha=0{,}8$, ambele oferte sînt acceptate în \NqbBoth\,\% din zile).
În operarea transfrontalieră, cea mai bună dintre rutele RO/BG/HU adaugă \NcrossGain\,\% la profitul obținut în ipoteza previziunii perfecte; cea mai frecventă rută transfrontalieră este \NtopRoute.
\textbf{Cifra relevantă pentru decizie:} în 2019--2024, strategia realistă a cîștigat \NqbYear\ EUR pe an, față de \NneedYear\ EUR necesari (VAN \NqbNPV\ EUR, perioadă de recuperare de \NqbPay\ ani) --- investiția nu este viabilă, ca și în lucrare.
În 2025--2026, cînd prețurile de la prînz au scăzut și mai mult din cauza producției solare (perechea fixă aleasă pe datele din 2024 este \NnhPairNew), strategia cîștigă \NqbYearNew\ EUR pe an (perioadă de recuperare de \NqbPayNew\ ani): rezultatul se apropie de prag, dar rămîne sub acesta în absența veniturilor din piața de echilibrare sau din plățile pentru capacitate.}

\section{\T{Discussion questions for managers}{Întrebări pentru discuție}}
\begin{enumerate}
\item \T{The vendor shows the perfect-foresight profit. By how much would you discount it, and why?}%
{Ofertantul prezintă profitul obținut în ipoteza previziunii perfecte. Cu cît l-ați ajusta și de ce?}
\item \T{The fixed pair chosen on the previous year moved from night$\rightarrow$evening to midday$\rightarrow$evening. What does that tell you about the value of co-locating a battery with your solar park?}%
{Perechea fixă de ore aleasă pe anul anterior s-a mutat de la noapte$\rightarrow$seară la prînz$\rightarrow$seară. Ce arată acest lucru despre valoarea amplasării unei baterii lîngă parcul fotovoltaic?}
\item \T{Which other revenues (balancing, frequency reserves, capacity payments, grid-fee savings) would you add to close the gap, and how certain are they?}%
{Ce alte venituri (din piața de echilibrare, din rezervele de reglaj al frecvenței, din plățile pentru capacitate, din economiile la tarifele de rețea) ați lua în calcul pentru a acoperi diferența și cît de sigure sînt acestea?}
\item \T{Why do safer bids ($\alpha=0.8$) earn more than aggressive ones? What does this say about trading on point forecasts?}%
{De ce cîștigă ofertele prudente ($\alpha=0{,}8$) mai mult decît cele agresive? Ce arată acest lucru despre tranzacționarea pe baza prognozelor punctuale?}
\item \T{At what CAPEX would the battery be viable on 2025--2026 prices? Is that a realistic price trend?}%
{La ce nivel al CAPEX ar deveni bateria viabilă la prețurile din 2025--2026? Este realistă o astfel de evoluție a prețurilor bateriilor?}
\end{enumerate}

\section{\T{Limitations of the classroom replication}{Limitele replicării}}
\T{(i) We leave out cross-border fees, congestion, intraday and balancing markets, and price impact.
(ii) One cycle per day and 1 MWh: real projects use 2--4 hour batteries and several revenue streams.
(iii) Perfect foresight is chosen ex post and is an upper bound; the fixed hours (chosen on the previous year) and the forecast-based rule are achievable.
(iv) No degradation cost and no grid tariffs (a few lei/MWh in Romania).
(v) Our forecasts use no fundamentals; the paper's DRF uses 37 predictors.}%
{(i) Nu luăm în calcul tarifele transfrontaliere, congestiile, piețele intrazilnice și de echilibrare și nici efectul ofertelor bateriei asupra prețului.
(ii) Un singur ciclu pe zi și o capacitate de 1 MWh: proiectele reale folosesc baterii cu durata de descărcare de 2--4 ore și mai multe surse de venit.
(iii) Previziunea perfectă este stabilită \emph{ex post} și reprezintă o limită superioară; orele fixe (alese pe anul anterior) și regula bazată pe prognoză pot fi aplicate în practică.
(iv) Nu includem costul degradării bateriei și nici tarifele de rețea (cîțiva lei/MWh în România).
(v) Prognozele noastre nu folosesc variabile fundamentale; modelul DRF din lucrare folosește 37 de variabile explicative.}

\newpage
\section*{\T{Case 6 at a glance}{Studiul de caz 6 pe scurt}}
\glance{\T{Does day-ahead arbitrage repay a 1 MW / 1 MWh battery when bids are set a day ahead?}{Își recuperează o baterie de 1 MW / 1 MWh investiția din arbitrajul pe PZU, dacă ofertele se depun cu o zi înainte?}}%
{\T{Hourly RO, BG, HU prices, 2019--Oct 2026; paper's window 2019--Sep 2024.}{Prețuri orare RO, BG, HU, 2019--oct.\ 2026; perioada din lucrare: 2019--sept.\ 2024.}}%
{\T{Random-forest probabilistic forecasts; quantile bids; fixed hours and perfect foresight as benchmarks; NPV.}{Prognoze probabilistice cu Random Forest; oferte la nivelul cuantilelor; orele fixe și previziunea perfectă ca repere; VAN.}}%
{\T{Realistic: \NqbYear\ EUR/yr (2019--2024), \NqbYearNew\ (2025--2026); needed \NneedYear.}{Strategia realistă: \NqbYear\ EUR/an (2019--2024), \NqbYearNew\ (2025--2026); necesar: \NneedYear.}}%
{\T{Not viable on arbitrage alone at 500k EUR CAPEX; stack revenues or wait for cheaper batteries.}{Neviabilă dacă veniturile provin exclusiv din arbitraj, la un CAPEX de 500.000 EUR; cumulați mai multe surse de venit sau așteptați scăderea prețului bateriilor.}}
\begin{center}
\includegraphics[width=0.66\linewidth]{fig_headline_\lang}
\src{\T{our computation on energy-charts.info prices, 2019--Oct 2026.}{calcule proprii pe prețurile energy-charts.info, 2019--oct.\ 2026.}}
\end{center}
\begin{minipage}[t]{0.32\linewidth}\centering{\color{MainBlue}\Huge\bfseries \NpfCum}\\ \small \T{EUR, perfect foresight 2019--2024 (paper 188,249)}{EUR, previziune perfectă, 2019--2024 (lucrarea: 188.249)}\end{minipage}\hfill
\begin{minipage}[t]{0.32\linewidth}\centering{\color{MainBlue}\Huge\bfseries \NqbPay}\\ \small \T{years payback, realistic strategy}{ani, perioada de recuperare pentru strategia realistă}\end{minipage}\hfill
\begin{minipage}[t]{0.32\linewidth}\centering{\color{IDAred}\Huge\bfseries \NqbNPV}\\ \small \T{EUR NPV at 500k CAPEX}{EUR, VAN la un CAPEX de 500.000 EUR}\end{minipage}
\end{document}
